% Tabla 4: Operacionalización de las variables.
% Etiqueta original: tab:operacionalizacion.
% Fuente: capitulos/3-metodologia.tex, lineas 64-86.
% Fragmento extraido automaticamente; requiere el preambulo de main.tex.

\begin{table}[htbp]
\centering\footnotesize
\caption{Operacionalización de las variables.}
\label{tab:operacionalizacion}
\begin{tabularx}{\textwidth}{@{}p{2.3cm} p{1.3cm} X X p{1.2cm}@{}}
\toprule
Variable & Tipo & Definición conceptual & Definición operacional & Unidad \\
\midrule
\multicolumn{5}{@{}l}{\emph{Factores variados}}\\
Método & Nominal & Ruta de conversión píxel--milímetro & A, \Aprime, B o C aplicados a los mismos puntos (Sección~\ref{sec:metodos}) & --- \\
Distancia al vidrio & Nominal & Profundidad $z$ del plano de simetría de la réplica & Posición del carrito: 5, 11, 17 o 23~cm; separación de la zona de B: 0, 6, 12 o 18~cm & cm \\
Rango de $z$ & Nominal & Conocimiento disponible sobre la profundidad & Toda la pecera (0--30~cm); $\pm 3$~cm; $\pm 1$~cm & cm \\
\multicolumn{5}{@{}l}{\emph{Resultados medidos}}\\
LT y PC estimadas & Razón & Longitud total y profundidad corporal de la réplica & Salida de Fishmeter por método, ecuación~\eqref{eq:medidas} & mm \\
LT y PC de referencia & Razón & Dimensión física de la réplica & Media de tres lecturas con calibre analógico, registradas al milímetro & mm \\
Error con signo, EA, EPA & Razón & Diferencia entre estimación y referencia & Ecuaciones~\eqref{eq:errores} & mm, mm, \% \\
Ancho relativo de $\IL$ & Razón & Extensión del intervalo por profundidad desconocida & $(\max_z L - \min_z L)/\mathrm{LT}_{\mathrm{ref}}$ & \% \\
Cobertura de $\IL$ & Razón & Fracción de capturas cuya referencia cae en $\IL$ & Conteo sobre las capturas disponibles & --- \\
\multicolumn{5}{@{}l}{\emph{Condiciones fijas}}\\
Cámara, pecera, agua, luz, marcadores, soporte, operador & --- & Condiciones del montaje & Registradas en las Secciones~\ref{sec:montaje} y~\ref{sec:procedimiento} & --- \\
\bottomrule
\end{tabularx}
\end{table}
